Sections~\ref{sec:ablation} and \ref{sec:profabl} treat profiling data as something the prompt either
carries or does not. A third option is to carry it \emph{interpreted}: to convert the metrics into
named remedies before the model sees them. Recent work argues strongly for this---KernelPro turns
Nsight Compute, SASS and Nsight Systems output into natural-language guidance through a set of
micro-profiling tools, and reports a $125\%$ higher speedup than handing the model raw metrics
\cite{gai2026kernelpro}. We evaluated the same idea, as a \texttt{diagnostics} prompt that adds a
bounded checklist of measured conditions and the transformation each one suggests, and reach a more
qualified conclusion: interpretation does not shift the average outcome so much as it widens the
distribution of outcomes.

\paragraph{Setup.}
Table~\ref{tab:guidvar} reports four arms on one node (GH200, idle, all arms built once and timed
round-robin in a single session): the unmodified program, the constrained \emph{baseline} prompt, the
relaxed \emph{structural} prompt, and \emph{diagnostics}, which is the structural prompt plus the
checklist. These are the tightest measurements in this paper---most repetition spreads are below $5\%$
and several are below $0.1\%$---so the differences below are not measurement artifacts.

\begin{table}[t]
\centering
\footnotesize
\setlength{\tabcolsep}{2pt}
\caption{Effect of converting profiling metrics into a prescriptive checklist (\emph{diag}), against
  the constrained (\emph{base}) and relaxed (\emph{struct}) prompts. Speedups are interleaved medians
  against the unmodified program on one GH200. \textbf{Bold} marks the cases discussed in the text.
  ``spread'' is the largest repetition spread across the four arms.}
\label{tab:guidvar}
\resizebox{\columnwidth}{!}{%
\begin{tabular}{@{}llrrrr@{}}
\toprule
\textbf{Benchmark} & \textbf{Model} & \textbf{base} & \textbf{struct} & \textbf{diag} & \textbf{spread} \\
\midrule
\texttt{chemv}                & Gemini & \textbf{7.835}$\times$ & 7.638$\times$ & \textbf{1.075}$\times$ & 1.9\% \\
\texttt{concat}               & GPT-5  & 1.031$\times$ & 0.838$\times$ & \textbf{0.349}$\times$ & 0.1\% \\
\texttt{boxfilter}            & GPT-5  & \textbf{10.280}$\times$ & 6.352$\times$ & 5.503$\times$ & 0.0\% \\
\midrule
\texttt{bincount}             & GPT-5  & 1.035$\times$ & 1.037$\times$ & \textbf{1.374}$\times$ & 5.2\% \\
\texttt{bincount}             & Gemini & 1.106$\times$ & 1.003$\times$ & \textbf{1.380}$\times$ & 4.4\% \\
\texttt{concat}               & Gemini & 1.001$\times$ & 1.797$\times$ & \textbf{1.798}$\times$ & 46.8\% \\
\texttt{boxfilter}            & Gemini & 8.691$\times$ & 8.930$\times$ & 9.371$\times$ & 42.6\% \\
\texttt{convolutionSeparable} & Gemini & 1.418$\times$ & 1.499$\times$ & 1.514$\times$ & 19.3\% \\
\midrule
\texttt{blockAccess}          & both   & $\approx 1.00\times$ & $\approx 1.00\times$ & $\approx 1.00\times$ & 0.2\% \\
\bottomrule
\end{tabular}%
}
\end{table}

\paragraph{The checklist produces our largest gains.}
On \texttt{bincount} it takes both models from roughly $1.04$--$1.11\times$ to $1.37$--$1.38\times$,
approaching our hand-tuned reference ($1.392\times$), and on \texttt{concat} under Gemini it is the
difference between no speedup and $1.80\times$. Where the measured condition is real and the named
remedy is the right one, telling the model what to do works.

\paragraph{It also produces our largest losses, on the same benchmarks.}
On \texttt{chemv} the unaided model had restructured the polyhedral-generated kernel into two
row-parallel kernels---one per triangle half, conjugation preserved on the upper one---a genuine
algorithmic rewrite worth $7.835\times$. Given the checklist, it kept the original loop nest and bolted
on the shared-memory staging the checklist named, collapsing to $1.075\times$. On \texttt{concat},
GPT-5 followed the instruction to process several elements per thread literally, to
one-thread-per-128-element line: a $128\times$ collapse in parallelism, measured at $0.349\times$ with a
$0.1\%$ spread. On \texttt{boxfilter} the same prompt costs GPT-5 nearly half of a $10.280\times$ result.

\paragraph{Interpretation narrows the search.}
A checklist of remedies is an anchor as much as a hint. When the listed transformation is the best
available one, anchoring is worth a large factor; when the model would have found something better
unaided, the anchor is what it finds instead. Averaged over our benchmarks the effect is close to
neutral, which is why a mean speedup---the statistic the literature reports---hides it. The dispersion
is the result: the same prompt component that produces $1.380\times$ on one benchmark produces
$0.349\times$ on another, both at sub-$5\%$ spread.

Two practical consequences follow. First, prescriptive guidance should be reported with its variance,
not only its mean; we would not have seen this in a geometric-mean summary. Second, the guidance should
be \emph{evidence rather than remedy}---state the measured condition (``$23\%$ occupancy, limited by
109 registers per thread'') and let the model decide what follows, instead of naming the
transformation. We did not get to test that variant, and it is the first thing we would run next.
